\documentclass[aspectratio=169,12pt]{beamer}

% ------------------------------------------------------------------
% Pacotes e tema
% ------------------------------------------------------------------
\usepackage[utf8]{inputenc}
\usepackage[T1]{fontenc}
\usepackage[brazil]{babel}
\usepackage{lmodern}

\usetheme{Madrid}
\usecolortheme{whale}
\useinnertheme{circles}

\setbeamertemplate{navigation symbols}{}
\setbeamertemplate{caption}[numbered]
\setbeamerfont{frametitle}{size=\large}

% ------------------------------------------------------------------
% Dados da capa
% ------------------------------------------------------------------
\title[Sentinela Cyber]{Diagnóstico Organizacional da Sentinela Cyber Consultoria}
\subtitle{Conceitos de Administração aplicados a uma consultoria em segurança da informação}
\author{Bruno Viana \and Felipe Koike \and Jônatas Brito}
\institute{Universidade Presbiteriana Mackenzie \\ Administração de Negócios -- Avaliação N2}
\date{2026}

\begin{document}

% ------------------------------------------------------------------
% Capa
% ------------------------------------------------------------------
\begin{frame}[plain]
  \titlepage
\end{frame}

% ------------------------------------------------------------------
% Agenda
% ------------------------------------------------------------------
\begin{frame}{Agenda}
  \begin{enumerate}
    \item Apresentação da empresa
    \item Estrutura organizacional
    \item Diagnóstico do ambiente: cultura e desafios
    \item Processo administrativo na prática
    \item Relacionamento interpessoal e soft skills
    \item Evolução do pensamento administrativo
    \item Reflexão crítica e recomendações
    \item Conclusão
  \end{enumerate}
\end{frame}

% ------------------------------------------------------------------
% Introdução
% ------------------------------------------------------------------
\begin{frame}{Contexto}
  \begin{block}{O problema}
    A transformação digital ampliou a exposição das organizações a ameaças
    cibernéticas, e o \textit{phishing} explora principalmente o comportamento
    humano.
  \end{block}
  \begin{block}{A proposta}
    Aplicar os conceitos de Administração à realidade de uma consultoria em
    segurança da informação focada no fator humano: a Sentinela Cyber.
  \end{block}
\end{frame}

% ------------------------------------------------------------------
% Apresentação da empresa
% ------------------------------------------------------------------
\begin{frame}{Apresentação da empresa}
  \textbf{Sentinela Cyber Consultoria} -- consultoria em segurança da
  informação centrada nas pessoas.

  \vspace{0.3cm}
  Fundada por três alunos do Mackenzie, com trajetórias complementares:
  \begin{itemize}
    \item Bruno Viana: 10 anos em auditoria;
    \item Felipe Koike: 9 anos como analista de dados sênior em banco;
    \item Jônatas Brito: 4 anos em infraestrutura e 3 anos como desenvolvedor.
  \end{itemize}

  \vspace{0.2cm}
  \begin{block}{Serviços}
    Treinamentos antiphishing, boas práticas no uso de ferramentas da internet,
    consultoria em políticas de segurança e simulações de engenharia social.
  \end{block}
\end{frame}

% ------------------------------------------------------------------
% Estrutura organizacional
% ------------------------------------------------------------------
\begin{frame}{Estrutura organizacional}
  Três áreas funcionais, cada uma com um time padronizado de 7 pessoas (21 no total):
  \begin{itemize}
    \item 1 Diretor (sócio fundador);
    \item 1 PM ou PO;
    \item 4 colaboradores;
    \item 1 estagiário.
  \end{itemize}

  \vspace{0.2cm}
  \begin{columns}[t]
    \column{0.33\textwidth}
      \begin{block}{Auditoria e Compliance}
        Bruno Viana
      \end{block}
    \column{0.33\textwidth}
      \begin{block}{Dados e Inteligência}
        Felipe Koike
      \end{block}
    \column{0.33\textwidth}
      \begin{block}{Tecnologia e Operações}
        Jônatas Brito
      \end{block}
  \end{columns}
\end{frame}

% ------------------------------------------------------------------
% Cultura
% ------------------------------------------------------------------
\begin{frame}{Cultura organizacional: 4 pilares}
  \begin{columns}[t]
    \column{0.5\textwidth}
      \begin{block}{Defesa Coletiva}
        A segurança é responsabilidade de todos.
      \end{block}
      \begin{block}{Confiança em Rede}
        A união entre colaboradores como ativo estratégico.
      \end{block}
    \column{0.5\textwidth}
      \begin{block}{Integridade Inegociável}
        Ética, transparência e sigilo como condição.
      \end{block}
      \begin{block}{Evolução Contínua}
        Aprendizado permanente diante de ameaças em mutação.
      \end{block}
  \end{columns}

  \vspace{0.2cm}
  Liderança participativa e clima de proximidade entre as equipes.
\end{frame}

% ------------------------------------------------------------------
% Desafios
% ------------------------------------------------------------------
\begin{frame}{Principais desafios}
  \begin{itemize}
    \item Dependência das pessoas-chave (os três sócios fundadores);
    \item Construção de credibilidade em um mercado novo;
    \item Atualização técnica constante frente a novas ameaças;
    \item Formalização de processos ainda informais;
    \item Mensuração de resultados intangíveis dos treinamentos.
  \end{itemize}
\end{frame}

% ------------------------------------------------------------------
% Processo administrativo
% ------------------------------------------------------------------
\begin{frame}{Processo administrativo na prática}
  \begin{columns}[t]
    \column{0.5\textwidth}
      \begin{block}{Planejamento}
        Definição do nicho, metas trimestrais e plano de ação por projeto.
      \end{block}
      \begin{block}{Organização}
        Divisão por áreas, padronização de metodologias e etapas de projeto.
      \end{block}
    \column{0.5\textwidth}
      \begin{block}{Direção}
        Liderança participativa, coordenação das equipes e engajamento.
      \end{block}
      \begin{block}{Controle}
        Indicadores: taxa de cliques, taxa de reporte e conclusão de treinos.
      \end{block}
  \end{columns}
\end{frame}

% ------------------------------------------------------------------
% Ferramentas de apoio à decisão
% ------------------------------------------------------------------
\begin{frame}{Ferramentas de apoio à decisão}
  \begin{itemize}
    \item Painéis (\textit{dashboards}) de indicadores de segurança e desempenho;
    \item Plataformas de simulação de phishing;
    \item Sistemas de gestão de projetos (ferramentas \textit{kanban});
    \item Matrizes de risco e análise SWOT.
  \end{itemize}
\end{frame}

% ------------------------------------------------------------------
% Soft skills
% ------------------------------------------------------------------
\begin{frame}{Relacionamento interpessoal e soft skills}
  Negócio técnico, mas dependente de competências comportamentais:
  \begin{itemize}
    \item Comunicação e didática na condução dos treinamentos;
    \item Empatia para abordar erros humanos sem constrangimento;
    \item Pensamento crítico e ético no acesso a informações sensíveis;
    \item Colaboração entre as diferentes especialidades.
  \end{itemize}

  \vspace{0.2cm}
  A combinação de \textit{hard skills} e \textit{soft skills} é o principal
  diferencial competitivo.
\end{frame}

% ------------------------------------------------------------------
% Evolução do pensamento administrativo
% ------------------------------------------------------------------
\begin{frame}{Evolução do pensamento administrativo}
  \begin{itemize}
    \item Administração Científica (Taylor): padronização e medição de resultados;
    \item Teoria Clássica (Fayol): funções administrativas e divisão de responsabilidades;
    \item Relações Humanas (Mayo): motivação, clima e valorização do fator humano;
    \item Escola Sistêmica: segurança como sistema de pessoas, processos e tecnologia.
  \end{itemize}
\end{frame}

\begin{frame}{Conceitos aplicados x superados}
  \begin{columns}[t]
    \column{0.5\textwidth}
      \begin{block}{Ainda aplicados}
        \begin{itemize}
          \item Padronização e métricas (Taylor);
          \item Funções administrativas (Fayol);
          \item Valorização do fator humano (Mayo).
        \end{itemize}
      \end{block}
    \column{0.5\textwidth}
      \begin{block}{Superados}
        \begin{itemize}
          \item Rigidez hierárquica;
          \item Visão mecanicista do trabalhador;
          \item Comando centralizado.
        \end{itemize}
      \end{block}
  \end{columns}
\end{frame}

% ------------------------------------------------------------------
% Reflexão crítica
% ------------------------------------------------------------------
\begin{frame}{Reflexão crítica e recomendações}
  \begin{block}{Teoria x prática}
    Base em princípios clássicos, mas dinâmica cotidiana humanística e
    sistêmica. Lacuna principal: baixa formalização de processos.
  \end{block}
  \begin{block}{Recomendações}
    \begin{itemize}
      \item Formalizar e documentar processos;
      \item Planejar a expansão das equipes;
      \item Adotar KPIs internos e capacitação contínua;
      \item Mensurar o valor percebido pelos clientes.
    \end{itemize}
  \end{block}
\end{frame}

% ------------------------------------------------------------------
% Conclusão
% ------------------------------------------------------------------
\begin{frame}{Conclusão}
  \begin{itemize}
    \item A gestão contemporânea é cumulativa: as escolas coexistem e se complementam;
    \item O diferencial da Sentinela Cyber está na complementaridade dos sócios e na cultura de quatro pilares;
    \item O desafio futuro é equilibrar flexibilidade e capital humano com a estruturação de processos.
  \end{itemize}
\end{frame}

\begin{frame}[plain]
  \centering
  \vfill
  {\Large Obrigado!}\\[0.4cm]
  {\normalsize Sentinela Cyber Consultoria}\\
  {\normalsize Bruno Viana \quad Felipe Koike \quad Jônatas Brito}
  \vfill
\end{frame}

\end{document}
